\subsection{PASSAGE FROM LIE ALGEBRA MORPHISMS TO LIE GROUP MORPHISMS}

Lemma 1. Let G be a simply connected† topological group, W a symmetric connected open neighbourhood of e, H a group and f a mapping of W³ into H such that

\[
f (x y z) = f (x) f (y) f (z)
\]

for \(x, y, z\) in W. There exists a morphism \(f'\) of G into H such that \(f' \mid \mathrm{W} = f \mid \mathrm{W}\) . For \((g, h) \in \mathrm{G} \times \mathrm{H}\) and U an open neighbourhood of \(e\) in W, let \(\mathrm{A}(g, h, \mathrm{U})\) be the set of \((gu, hf(u)) \in \mathrm{G} \times \mathrm{H}\) for \(u \in \mathrm{U}\) . Then \((g, h) \in \mathrm{A}(g, h, \mathrm{U})\) and \(\mathrm{A}(g, h, \mathrm{U}_1) \cap \mathrm{A}(g, h, \mathrm{U}_2) = \mathrm{A}(g, h, \mathrm{U}_1 \cap \mathrm{U}_2)\) . Let \((s, t) \in \mathrm{A}(g, h, \mathrm{U})\) ; then \(s = gu\) and \(t = hf(u)\) for \(u \in \mathrm{U}\) ; there exists an open neighbourhood \(\mathrm{U}'\) of \(e\) in W such that \(u\mathrm{U}' \in \mathrm{U}\) ; then, for \(u' \in \mathrm{U}'\) ,

\[
(s u ^ {\prime}, t f (u ^ {\prime})) = (g u u ^ {\prime}, h f (u u ^ {\prime})) \in \mathrm{A} (g, h, \mathrm{U})
\]

and hence \(\mathrm{A}(s,t,\mathrm{U}^{\prime})\subset \mathrm{A}(g,h,\mathrm{U})\) . It then follows that the \(\mathrm{A}(g,h,\mathrm{U})\) form the base of a topology on \(\mathbf{G}\times \mathbf{H}\) . We shall denote by \(\mathbf{Y}\) the set \(\mathbf{G}\times \mathbf{H}\) with this topology and denote by \(p\) the canonical projection of \(\mathbf{Y}\) onto \(\mathbf{G}\) , which is open. The restriction of \(p\) to \(\mathrm{A}(g,h,\mathrm{U})\) is a homeomorphism of \(\mathrm{A}(g,h,\mathrm{U})\) onto \(g\mathrm{U}\) . Hence \((\mathrm{Y},p)\) is a covering space of \(\mathbf{G}\) . Let \(\mathrm{Y}_0\) be the subgroup of \(\mathrm{Y}\) generated by \(\mathrm{A}(e,e,\mathrm{W})\) and let \(\mathfrak{B}\) be the set of \(\mathrm{A}(e,e,\mathrm{U})\) . Clearly \(\mathfrak{B}\) satisfies conditions \((\mathrm{GV}_{\mathrm{I}}^{\prime})\) and \((\mathrm{GV}_{\mathrm{II}}^{\prime})\) of General Topology, Chapter III, § 1, no. 2. The set \(\mathrm{Y}_0^{\prime}\) of \(y\in \mathrm{Y}_0\) such that the mappings \(z\mapsto yzy^{-1}\) and \(z\mapsto y^{-1}zy\) of \(\mathrm{Y}_0\) into \(\mathrm{Y}_0\) are continuous at \((e,e)\) is a subgroup of \(\mathrm{Y}_0\) . Let \(w\in W\) . The mapping \(w'\mapsto ww'w^{-1}\) of \(\mathrm{W}\) into \(\mathrm{G}\) is continuous and hence the mapping

\[
\left(w ^ {\prime}, f \left(w ^ {\prime}\right)\right) \mapsto \left(w w ^ {\prime} w ^ {- 1}, f \left(w w ^ {\prime} w ^ {- 1}\right)\right)
\]

of \(\mathrm{A}(e,e,\mathrm{W})\) into \(\mathrm{Y}_0\) is continuous at \((e,e)\). Now

\[
f \left(w w ^ {\prime} w ^ {- 1}\right) = f (w) f \left(w ^ {\prime}\right) f \left(w ^ {- 1}\right)
\]

\[
\left(w w ^ {\prime} w ^ {- 1}, f \left(w w ^ {\prime} w ^ {- 1}\right)\right) = (w, f (w)) \left(w ^ {\prime}, f \left(w ^ {\prime}\right)\right) \left(w, f (w)\right) ^ {- 1}
\]

As \(w^{-1} \in W\) , we see that \((w, f(w)) \in Y_0'\) . Thus, \(A(e, e, W) \subset Y_0'\) , so that \(Y_0' = Y_0\) . The group \(Y_0\) , with the filter base \(\mathfrak{B}\) , therefore satisfies condition \((GV_{\mathrm{III}}')\) of General Topology, Chapter III, § 1, no. 2. As

\[
(g, h). \mathrm{A} (e, e, \mathrm{U}) = \mathrm{A} (g, h, \mathrm{U}),
\]

\(Y_{0}\) is a topological group which is connected since \(A(e, e, W)\) is connected. Then \(p(Y_{0})\) is an open subgroup of G, whence \(p(Y_{0}) = G\) since G is connected. The kernel of \(p \mid Y_{0}\) is discrete. As G is simply connected, \(p \mid Y_{0}\) is a homeomorphism of \(Y_{0}\) onto G. Therefore \(Y_{0}\) is the graph of a morphism \(f'\) of G into H. For \(g \in W\) , \((g, f(g)) \in A(e, e, W) \subset Y_{0}\) , whence \(f(g) = f'(g)\) .

THEOREM 1. Let G and H be Lie groups and h a continuous morphism of L(G) into L(H). Suppose that G is simply connected. Then there exists one and only one Lie group morphism \(\phi\) of G into H such that \(h = \mathrm{L}(\phi)\) .

The existence of \(\phi\) follows from Lemma 1 and §4, no. 1, Theorem 1 (i). The uniqueness of \(\phi\) follows from §4, no. 1, Theorem 1 (ii) and the fact that \(G\) is connected.

COROLLARY. Let G be a finite-dimensional simply connected Lie group. There exists a finite-dimensional analytic linear representation of G whose kernel is discrete.

There exist (Chapter I, § 7, Theorem 2) a finite-dimensional vector space E and an injective morphism \(h\) of \(L(G)\) into the Lie algebra \(\operatorname{End}(E)\) . By Theorem 1, there exists a morphism \(\phi\) of G into \(\mathbf{GL}(E)\) such that \(L(\phi) = h\) . Hence \(\phi\) is an immersion and therefore its kernel is discrete.

Remarks. (1) There exist finite-dimensional simply connected Lie groups which have no finite-dimensional injective analytic linear representation (Exercise 2).

\begin{enumerate}
\def\labelenumi{(\arabic{enumi})}
\setcounter{enumi}{1}
\tightlist
\item
  There exist finite-dimensional connected Lie groups G such that every finite-dimensional analytic linear representation of G has non-discrete kernel (Exercises 3 and 4).
\end{enumerate}

